% ===== sec/03_generator_verifier_architecture.tex =====
\section{Generator--Verifier Architecture}
\label{sec:architecture}

The architecture has two principal agents: a Goal-Driven generator
$\mathcal{G}$ and a Data-Driven verifier $\mathcal{V}$.
Theorem~\ref{thm:closure} guarantees their composition preserves the
generator's class; this section makes the construction concrete.

\subsection{The two principal agents}

$\mathcal{G}$ is modelled by an idealised unrestricted generator,
realised concretely as a bounded LLM loop.  Its step relation
$\delta_\mathcal{G}$ samples a continuation from the LLM, parses
tool-call markup, and emits the corresponding operation string.
The model used is \texttt{gemini-2.5-pro} for deliberation steps and
\texttt{gemini-2.5-flash-lite} for housekeeping.

 $\mathcal{G}$ receives at start-up the start-language alphabet
$\Sigma_0$ (the columns of
\texttt{trace\_language\allowbreak \_original.csv}), an empty workspace,
and the Kaggle CLI configuration.  It is \emph{not} given the goal label,
leaderboard specification, or any prior solution.

$\mathcal{V}$ is the Aho--Corasick DFA $M_K$ over a finite keyword
set $K$---the union of the \texttt{unit\_test\_runner} and
\texttt{tools\_index} columns of the start language, plus column-name
and dtype keywords from the task schema.  $\mathcal{V}$ has no LLM
component: its decision is
\[
h_\mathcal{V}(\tau) \;=\;
\begin{cases}
\mathsf{accept} & \text{if every required}\\
                 & \text{keyword class is matched} \\
\mathsf{reject} & \text{otherwise.}
\end{cases}
\]

\begin{proposition}[Verifier as DFA]
\label{prop:verifier-dfa}
$\mathcal{V}$ is realisable as a DFA with $O(\|K\|)$ states.  For
any $\tau\in\Sigma^*$, the verdict is computed in $O(|\tau|)$ time.
\end{proposition}

\noindent We choose Aho--Corasick over an LLM judge for three
reasons: (i) class containment---a DFA is Type-3 and gives $O(n)$
verifier cost by Corollary~\ref{cor:budget}; (ii) auditability---the
keyword set $K$ is finite, human-readable, and version-controlled;
(iii) robustness---a DFA cannot be ``talked into'' accepting a
malformed trace.

\subsection{The binding}

Algorithm~\ref{alg:enum} realises the binding
$\mathcal{G}\,\Vert\,\mathcal{V}$.

\begin{figure}[ht]
\noindent\rule{\columnwidth}{0.4pt}
\textbf{Algorithm 1.} \textsc{RecursivelyEnumerate}$(\mathcal{G},
\mathcal{V}, x)$
\par\noindent\rule{\columnwidth}{0.4pt}
\begin{algorithmic}[1]
\REQUIRE $\mathcal{G}$, $\mathcal{V}$, input $x$.
\ENSURE Verified trace $\tau$. \STATE
$q\gets q_0^\mathcal{G}$;\quad $\tau\gets\varepsilon$;\quad
$s\gets q_0^\mathcal{V}$.
\LOOP
  \STATE $(q', w) \gets \delta_\mathcal{G}(q, x)$.
  \STATE $s' \gets \delta_\mathcal{V}^*(s, w)$.
  \IF{$s'$ is a reject state}
    \STATE feed witness back to $\mathcal{G}$ as percept.
  \ELSE
    \STATE $\tau \gets \tau\cdot w$;\quad $q\gets q'$;\quad
    $s\gets s'$.
    \IF{$q\in F_\mathcal{G}$ and $s\in F_\mathcal{V}$}
      \RETURN $\tau$.
    \ENDIF
  \ENDIF
\ENDLOOP
\end{algorithmic}
\rule{\columnwidth}{0.4pt}
\caption{Recursive enumeration.}
\label{alg:enum}
\end{figure}

\subsection{Projected sub-agent views}

Intermediate sub-agent roles are projections of the validated trace
language, not independently designed machines.

\begin{definition}[Projection]
\label{def:projection}
Let $L(\mathcal{A})\subseteq \Sigma_\mathcal{A}^*$ and
$\Sigma' \subseteq \Sigma_\mathcal{A}$.  The \emph{projection} of
$L(\mathcal{A})$ onto $\Sigma'$ is
\[
\pi_{\Sigma'}(L(\mathcal{A})) \;=\;
\{\,\pi_{\Sigma'}(\tau) \,:\, \tau\in L(\mathcal{A})\,\}
\]
where $\pi_{\Sigma'}$ deletes every symbol not in $\Sigma'$.
\end{definition}

\begin{theorem}[Derived sub-agents]
\label{thm:derived}
For each intermediate role (planner, reviewer, task orchestrator,
etc.), its projected trace language is no greater than $\mathcal{L}_0$
because projection cannot raise the generator envelope's class.
Sharper class labels (Type-2?, indexed?, Type-1?) are modelling
hypotheses, not consequences of the projection lemma.
\end{theorem}

\subsection{Chomsky classification}

All non-Type-3 class assignments are modelling hypotheses; only the
verifier $\mathcal{V}$ has an explicit DFA recogniser.
Table~\ref{tab:class-signature} summarises the empirical
class-hypothesis signature.  In the realised system, Type-3 agents
(\texttt{tools\_index}, \texttt{unit\_test\_runner}) are the
verifier itself.  Type-2 hypotheses (\texttt{planner\_agent},
\texttt{reviewer\_agent}) correspond to stack-discipline traces.
The \texttt{sop\_orchestrator} exhibits a finite-index repair loop
suggestive of Aho's indexed grammars~\cite{aho1968indexed}.
Type-1 hypotheses (\texttt{task\_orchestrator},
\texttt{kernel\_pipeline}) model linear-bounded workspace growth.
Type-0 hypotheses (\texttt{research\_agent},
\texttt{developer\_agent}, \texttt{mlebench\_solver},
\texttt{kaggle\_trainer}) are idealised recursive-enumeration
envelopes.

\begin{table}[ht]
\centering
\caption{Empirical class-hypothesis signature.}
\label{tab:class-signature}
\resizebox{\columnwidth}{!}{%
\begin{tabular}{l|l|l}
\hline
\textbf{Sub-agent} & \textbf{Abstraction} & \textbf{Recogniser} \\ \hline
$\mathcal{V}$ (verifier)     & Type-3   & DFA \\
\texttt{planner\_agent}      & Type-2?  & PDA hypothesis \\
\texttt{reviewer\_agent}     & Type-2?  & PDA hypothesis \\
\texttt{sop\_orchestrator}   & Indexed? & Indexed-grammar hypothesis \\
\texttt{task\_orchestrator}  & Type-1?  & LBA hypothesis \\
\texttt{kernel\_pipeline}    & Type-1?  & LBA hypothesis \\
$\mathcal{G}$ (generator)    & Type-0?  & TM hypothesis \\
\hline
\end{tabular}%
}
\end{table}

\subsection{Recursive enumeration outcome}

When Algorithm~\ref{alg:enum} runs on the \textsc{Spaceship Titanic}
input, $\Sigma_0$ extends to $\Sigma_1$ and a thirteenth column
appears---the \texttt{kaggle\_trainer}---whose symbols were absent
from the start language.  Its emergence is forced by the verifier
rejecting any trace that omits a training step.  The empirical CSV
is a fixed point of the binding:
$\Sigma_1 = \Sigma_0 \cup \mathrm{Range}(\delta_\mathcal{G})\big|_{\mathcal{V}\text{-accepting}}$.
